\documentclass[11pt,a4paper]{article}
\usepackage[T1]{fontenc}
\usepackage{lmodern,amsmath,amssymb,amsthm,booktabs,microtype}
\usepackage[margin=24mm]{geometry}
\usepackage[colorlinks=true,urlcolor=blue,linkcolor=blue]{hyperref}
\newcommand{\Js}{\mathcal J_s}
\newcommand{\av}[1]{\tau\!\left(#1\right)}
\newcommand{\W}{\mathcal W}
\newcommand{\DD}{\mathfrak D}
\newcommand{\Rh}{\widehat R}
\newcommand{\sig}{\sigma}
\DeclareMathOperator{\tr}{tr}
\title{Mixed quadratic and cubic derivative singlets\\A local reduction to the multiprobe chord law}
\author{Alan Day \and Chengheng\thanks{AI research assistant (OpenAI).}}
\date{Local research note --- 6 November 2026}
\begin{document}\maketitle
\begin{abstract}
The quadratic and cubic derivative singlets constructed in the preceding
notes obey the same joint fixed-degree moment rule as two independent
probe species of lengths two and three. The microscopic operators share
the Hamiltonian couplings; the reduction follows from cancellations of
marked Wick diagrams. We classify every local component into fourteen
templates and give an exact polynomial certificate. This supplies the
missing interface to the existing multiprobe four-point function. We
evaluate the mixed out-of-time-order correlator, restore its mutual
crossing factor $q^6$, and obtain its short-time commutator growth.
\end{abstract}
\section{Joint law and physical target}
Use standard CAR, the complex Gaussian number-preserving $p$-body
Hamiltonian including overlapping index sets, and the annealed normalized
trace $\tau=\mathbb E\tr$ at $\beta=\mu_{\rm chem}=0$. With $\hbar=1$, take
\[
 N,p\longrightarrow\infty,\quad\lambda=p^2/N>0,\qquad
 q=e^{-4\lambda},\qquad \sig^2=\Js^2/\lambda.
\]
Let $h=H/\sig$, $X=Y/\sqrt W$ and $\Rh=R/\sqrt{V_R}$ be the normalized
Hamiltonian, quadratic singlet, and cubic residual defined on page~2.
Their joint limiting moments are
\begin{equation}
 \boxed{\av{A_1\cdots A_n}
 =\sum_{\substack{\pi\text{ pairs identical species}}}
       \prod_{\substack{\{e,f\}\subset\pi\\ e\text{ crosses }f}}
       q^{d_e d_f},\qquad (d_h,d_X,d_{\Rh})=(1,2,3).}\label{law}
\end{equation}
In particular,
\begin{equation}
 \av{X\Rh^2}=0,\quad \av{X^2\Rh^2}=1,\quad
 \av{X\Rh X\Rh}=q^6,\quad \av{\Rh^4}=2+q^9.\label{moments}
\end{equation}
Thus distinct singlets already have a nonzero equal-time commutator:
$\av{[X,\Rh]^\dagger[X,\Rh]}=2(1-q^6)$.
The established multiprobe framework \cite{bint,bnr} now applies to these
specific composite insertions. Our task is the reduction to that framework,
not a replacement four-point kernel.

\newpage
\section{Operators and local marked weights}
Put $L=[H,\cdot]$, $\mathsf J_H B=\{H,B\}/2$ and
\[
 A=\sum_i(Lc_i)^\dagger Lc_i,\quad
 M_3=\sum_i c_i^\dagger L^3c_i,\quad T_3=(M_3+M_3^\dagger)/2.
\]
The preceding construction \cite{local} uses ordinary Gaussian coupling
chaos $\W_j$, not energy normal ordering. At finite $N$,
$Y_N=A-a_0-c_N(H^2-C)$ subtracts the scalar expectation and its quadratic
Hamiltonian projection, with the coefficients defined sector by sector.
Its DS limit is $Y$, and
\[
 c_N\longrightarrow\frac{2\lambda(1-q)}{1+q},\qquad
 W=4\Js^4\left[\lambda^{-1}+\frac{4q}{1+q}\right].
\]
For $(T,S,Z)=\W_3(T_3,\mathsf J_H A,H^3)$ set
\begin{align}
 R&=T-aS-bZ,&a&=-\frac{(1-q)^2}{1+q^2},\\
 b&=\frac{2\lambda(1-q)(1+q^3)}{(1+q^2)(1+q+q^2)},&
 D&=(1+q^2)(1+q+q^2).
\end{align}
Its variance is $V_R=\sig^6G_R/D^2$, where
\begin{align}
 G_X&=(1+q)^2W/\sig^4=4\lambda(1+q)^2+16\lambda^2q(1+q),\\
 G_R&=4\lambda(1-q^2)(1+q^2)(1+q+q^2)^2\notag\\
 &\quad+16\lambda^2q^2(1+q)(1+q^2)(1+q+q^2)(2q^2-q+2).
\end{align}
The symbol \texttt{X} in the raw polynomial code denotes $Y$, with the
denominator $1+q$ cleared; \texttt{R} has denominator $D$ cleared.
Every underlying Hamiltonian slot contributes $\sig$.

Expand a composite into consecutive slots. A quadratic vertex has one
possible mark, $12$; a cubic vertex has marks $12,23,13$. After clearing
denominators their weights are
\begin{center}
\begin{tabular}{ccl}
\toprule
vertex&mark coefficients&unmarked coefficient\\\midrule
$Y$&$1+q$&$-2\lambda(1-q)$\\
$R$&$-q(1+q+q^2),\ -q(1+q+q^2),\ D$&$-2\lambda(1-q)(1+q^3)$\\
\bottomrule
\end{tabular}
\end{center}
There is at most one mark per vertex. Internal Gaussian pairings within
each composite block are deleted before the DS limit. The exact identity
$A=(2p-\DD)H^2$, $\DD B=QB-\sum_i c_i^\dagger Bc_i$, and its three-slot
counterpart produce these marks, including the nonadjacent $13$ choice.

\newpage
\section{Why fourteen local templates suffice}
For an underlying pairing $\pi$, let $c(\pi)$ count ordinary chord
crossings. The previously established fixed-word determinant expansion
\cite{local} gives, for selected marks $j=(a,b)$,
\begin{equation}
 q^{c(\pi)}\exp\!\left[\lambda\sum_j\alpha_j z_j+
                 \lambda\sum_{j<k}\delta_{jk}z_jz_k\right],\qquad
 \alpha_j=2-4\chi_j.\label{marks}
\end{equation}
Here $\chi_j$ tests whether the two marked chords cross. For
$j=(a,b)$ and $k=(c,d)$, $\delta_{jk}=4$ precisely when
$\pi(a)=d,\pi(b)=c$; otherwise it is zero. Extract each selected mark
once and insert the coefficients on page~2. The exponential has no
higher connected mark term: triple set overlaps vanish in the fixed-word
DS limit. Fixed-order overlap exponential bounds justify the derivatives.

\paragraph{Component classification.}
Make a graph whose vertices are the composite blocks, joining two blocks
when at least two Hamiltonian chords connect them. A vertex has two or
three slots, so it has at most one such neighbor. Every component is a
singleton or a pair. Ordinary one-slot $H$ vertices remain spectators.
Fix the component's external endpoints and sum over all assignments to
its internal slots. All crossings with chords having no endpoint in the
component factor out: the component slots are consecutive blocks.

An outside mark cannot change under these assignments. If both of its
chords ended in one component vertex, it would itself be a double-bridge
neighbor, already in that component. If the partners end in distinct
vertices, their cyclic block order is fixed. A delta term connecting the
component to an outside mark would likewise require a double bridge.
Thus all outside marked factors are constant during the local sum.
Crossings among the component's own incident chords are retained in the
local template, including its residual external endpoints.

\begin{center}
\begin{tabular}{lrrl}
\toprule
local component&templates&local pairings&sum\\\midrule
singleton $Y$&1&2&0\\
singleton $R$&1&6&0\\
$YY$, two bridges&1&2&$G_X$\\
$YR$ or $RY$, two bridges&4&24&0\\
$RR$, two bridges&6&108&0\\
$RR$, three bridges&1&6&$G_R$\\\bottomrule
\end{tabular}
\end{center}
For two blocks, external endpoints can lie on either separating arc.
The six partial $RR$ templates are two endpoint color orders times three
arc splits. These possibilities exhaust the topology. Exact determinant
polynomials give the displayed sums for independent symbols $q,\lambda$.
Every incomplete component cancels. The survivors pair whole equal-species
blocks. Compressing such pairs gives $d_e d_f$ elementary crossings,
and dividing by $G_X,G_R$ proves \eqref{law} for every fixed word.

\newpage
\section{Real-time mixed interaction observable}
Define $X(t)=e^{itH}Xe^{-itH}$ and
\begin{equation}
 F_{23}(t)=\av{X(t)\Rh X(t)\Rh},\qquad
 D_{23}(t)=\av{[X(t),\Rh]^\dagger[X(t),\Rh]}.
\end{equation}
The uncrossed mixed pair is $\av{X(t)^2\Rh^2}=1$. Trace cyclicity and
Hermiticity therefore give $D_{23}(t)=2[1-\operatorname{Re}F_{23}(t)]$.
For distinct species there is only the single allowed crossed pairing
in $F_{23}$, rather than the three channels of a same-species correlator.

Here is a positive-series form of BINT Eq.~(4.10) \cite{bint}, in symmetric
Jacobi normalization. Write $a=q^2$, $b=q^3$, $Q_{23}=q^6$ and
\begin{align}
 J_{n-1,n}&=\sig\sqrt{\frac{1-q^n}{1-q}},&U(t)&=e^{-itJ},&v_n&=U_{n0}(t),\\
 c_i&=\frac{(b^2;q)_i}{(q;q)_i}a^i,&
 V_{li}&=\sqrt{c_i}\,b^l
       \sqrt{\frac{(q;q)_{l+i}}{(q;q)_l}}v_{l+i}.
\end{align}
With $(x;q)_i=\prod_{j=0}^{i-1}(1-xq^j)$,
\begin{equation}
 \boxed{F_{23}(t)=Q_{23}\sum_{n,i\ge0}a^n
                    |(U(t)^\dagger V)_{ni}|^2.}\label{series}
\end{equation}
BINT explicitly suppresses the overall mutual-probe factor in that
section. Restoring $Q_{23}$ is essential: $F_{23}(0)=q^6$, not one.
BNR Section~8.1 \cite{bnr} supplies the neutral complex-model interface.
In \eqref{series}, the four alternating real-time propagators have
combined into a modulus square; this is not an extra approximation.

Direct compressed chord enumeration gives
\begin{align}
 F_{23}(t)&=q^6-2q^6(1-q^2)(1-q^3)(\sig t)^2
                +c_4(q)(\sig t)^4+O(t^6),\\
 c_4(q)&=\frac{q^6(q-1)^3(q+1)(q^2+q+1)}6\notag\\
 &\qquad\times(9q^5+12q^4+9q^3-3q^2-12q-11),\\
 D_{23}(t)&=2(1-q^6)+4\sig^2q^6(1-q^2)(1-q^3)t^2+O(t^4).
\end{align}
In particular $D_{23}(0)=48\lambda+O(\lambda^2)$, and the coefficient
of $t^2$ is $384\lambda\Js^2+O(\lambda^2)$. The noncommutativity and its
initial growth both survive at the first deformation order. A
species-preserving Wick law does not make these ordinary commuting
Gaussian variables.

\newpage
\section{Verification and numerical results}
The initial mixed enumeration covers twelve words, 15,392 pairings and
433,964 local-minor checks, including
$X\Rh^2$, $X\Rh X\Rh$, $X^2\Rh^2$, $\Rh^4$, and words with two external
Hamiltonians. All differences from \eqref{law} are exact zero polynomials.
The separate fourteen-template certificate uses the older determinant
implementation, not the new geometric weight evaluator. It shares the
pairing generator. The factorization argument on page~3 is what extends
these finite checks to arbitrary fixed degree.

An independent finite-CAR computation tests all choices of local mark
or no mark on 64 mixed words at $(N,p)=(3,1),(4,2)$. All 6,656 normalized
matrix traces equal the finite determinant coefficients exactly. Among
them 1,031 are nonzero, including two negative values. Coherent repeated
index sets are included to avoid an all-zero sparse test. These checks
validate the mark algebra, not finite-size convergence rates.

The real-time code uses $\Js=1$. It independently sums all four-arc
Hamiltonian insertions through order $t^6$ and compares with
\eqref{series} at $t=.01$, swaps the two probe labels, and enlarges each
Jacobi cutoff by 16. Results at $t=1$ are
\begin{center}
\begin{tabular}{rrrr}
\toprule
$\lambda$&cutoff&$F_{23}(1)$&$D_{23}(1)$\\\midrule
.20&52&$4.26855521\,10^{-6}$&1.99999146\\
.10&90&$5.71877829\,10^{-5}$&1.99988562\\
.05&167&$5.68948591\,10^{-4}$&1.99886210\\\bottomrule
\end{tabular}
\end{center}
These are finite-$\lambda$ values, not estimates of a first-order
small-$\lambda$ coefficient.
Across all fifteen points, the largest cutoff and species-swap changes
are $6.39\,10^{-16}$ and $3.34\,10^{-16}$, respectively. The small-time
Taylor comparison differs by at most $1.27\,10^{-15}$.

\paragraph{Analytic truncation control.}
For $b=q^{d_2}$, put
\[
 C_b=\prod_{j=1}^{2d_2-1}(1-q^j)^{-1},\quad K=C_b/(1-a),\quad
 R_H=2\sig/\sqrt{1-q}.
\]
Then $c_i\le C_ba^i$, the map $v\mapsto V$ has norm at most $\sqrt K$,
and its rows $l\ge M$ have norm at most $\sqrt K b^M$. A cutoff of size
$L$ agrees with powers of the infinite Jacobi matrix up to the first
path reaching the boundary. Define $T_k(z)=\sum_{j\ge k}z^j/j!$ and
\[
 \delta=2T_L(R_H|t|)+2T_{L-M}(R_H|t|)+2b^M.
\]
Splitting low and high source rows, unitarity gives
$|F_{23}-F_{23}^{(L)}|\le2Q_{23}K\delta$.
We choose the larger cutoff needed for the two species orders; the
reported analytic bounds are below $10^{-11}$ for $|\Js t|\le1$.
This bounds truncation in exact arithmetic. Floating-point accuracy is
assessed by the separate comparisons, not certified by interval arithmetic.

\newpage
\appendix
\section{Fixed times, limits, and the next calculation}
Equation \eqref{law} is a fixed-degree DS statement. It also determines
fixed real-time correlators. Here is a moment-based justification that
does not assume a uniform finite-$N$ operator norm. For a finite real-time
word, approximate each unitary by its degree-$k$ Taylor polynomial.
The scalar spectral theorem bounds its $L^r(\tau_N)$ remainder by
$|t|^{k+1}\|H^{k+1}\|_r/(k+1)!$ for even integer $r$. Noncommutative
H\"older bounds the rest of the word by fixed moments of the probes.
For each fixed $k$, all these moments converge by \eqref{law}. The
limiting Hamiltonian has the compact $q$-Hermite support
$[-R_H,R_H]$, so the limiting remainder is bounded by a constant times
$(|t|R_H)^{k+1}/(k+1)!$ (and fixed polynomial factors from the other
Taylor approximants). It vanishes as $k\to\infty$. Thus DS is taken
first, followed by polynomial degree. No growing-time or uniform
$\lambda\to0$ estimate is used here.

The local determinant limit uses fixed numbers of slots and marks.
At $\mu_{\rm chem}=0$, the central charge window and its exponentially
small complement are handled as in the preceding notes. The two-body
overlap exponent and its mark derivatives are uniform in that window.
Gaussian chaos subtraction removes internal pairings before the limit.
The remainder coefficients in $Y_N$ converge in the fixed moment bounds;
replacing them by their limiting values therefore preserves the theorem.

\paragraph{Next constructive step.}
The saved exact coefficients through $t^6$ give, coefficient by coefficient,
\[
 \frac{D_{23}(t)}\lambda:
 \quad48+384(\Js t)^2+512(\Js t)^4+
                  \frac{4096}{15}(\Js t)^6+\cdots
 \qquad(\lambda\to0).
\]
These are the first four terms of $48\cosh(4\Js t)$. This suggests a
concrete next target: use the existing unequal-probe first-order
four-point expansion to determine the entire fixed-time coefficient,
then match it to the two-species Rindler construction. Finite Taylor
agreement alone is not the derivation of that function.

The present theorem covers $H,X,\Rh$. It does not yet include arbitrary
high-point insertions of the linear rank-filter residual $F$. The
time-dependent mixed commutator is microscopic interaction data; selecting
a bulk mass, local vertex, or spatial dictionary remains a separate
physical construction. All statements here use the ensemble-averaged
infinite-temperature state and the specified order of limits.

\begin{thebibliography}{9}
\bibitem{bint} M. Berkooz, M. Isachenkov, V. Narovlansky and G. Torrents,
\emph{Towards a full solution of the large N double-scaled SYK model},
\href{https://arxiv.org/abs/1811.02584}{arXiv:1811.02584v2},
Sections 4.1--4.2, especially Eq.~(4.10) and footnote~8.
\bibitem{bnr} M. Berkooz, V. Narovlansky and H. Raj,
\emph{Complex Sachdev-Ye-Kitaev model in the double scaling limit},
\href{https://arxiv.org/abs/2006.13983}{arXiv:2006.13983v2},
Section 8.1, Eqs.~(8.1)--(8.7).
\bibitem{local} A. Day and Chengheng, local project notes:
\emph{Singlet multimark}, \emph{Cubic singlet}, \emph{Cubic kernel},
and \emph{Filtered spectrum} (November 2026).
Exact scripts and certificates: \texttt{results/susskind-mixed-cubic}.
\end{thebibliography}
\end{document}
